\documentclass[10pt,a4paper,titlepage,twoside,openright]{report}
%packages
\usepackage{makeidx}
\usepackage[portuges]{babel}
\usepackage{fancyhdr}
\usepackage{lipsum} % for dummy text
\usepackage{enumitem}
\setlist{nosep} % or \setlist{noitemsep} to leave space around whole list
% Figures
\usepackage[usenames,dvipsnames]{color}
\usepackage{graphicx}
\usepackage{caption}
\usepackage{multirow}
\usepackage{framed}
% Required for custom colors
\usepackage[usenames,dvipsnames]{xcolor}
% hyperlinks
\usepackage{url}
\usepackage{hyperref}
% Math
\usepackage{amssymb,amsbsy,amsmath}
\usepackage{amsthm} %para poder definir newtheorem
\usepackage{amsfonts} %para poder fazer o R, Z etc \mathbb{R}	\mathbb{Z} 

%estilo de página
\pagestyle{fancy}
			\lhead{Trabalho Prático de Matemática Discreta}
			\chead{ }
			\rhead{	\includegraphics[height=0.45cm]{Figuras/PPorto-peq}}
			\lfoot{\hrule }
			\cfoot{}
			\rfoot{\hrule \vspace*{5pt} \thepage}
	
%\renewcommand{\chaptermark}[1]{olaallalal}


\linespread{1.5}    							% espaçamento entre linhas
\setlength{\parindent}{0pt} 					% espaço horizontal aquando a mudança de paragrafo
\setlength{\parskip}{1ex plus 0.5ex minus 0.2ex}% espaço vertical entre parágrafos
\setlength{\headheight}{17pt}
\newcommand{\Lower}[1]{\smash {\lower 1.5ex \hbox{#1}}}

\usepackage{titlesec}
\titleformat
{\chapter} % command
[display] % shape
{\bfseries\Huge%\itshape
} % format
{%Story No. \ \thechapter
} % label
{0.5ex} % sep
{
%    \rule{\textwidth}{1pt}
    \vspace{1ex}
    %\centering
} % before-code
[
%\vspace{-0.5ex}%
%\rule{\textwidth}{0.3pt}
] % after-code
%__________________________________________________________________________________________________________
\begin{document}



\begin{titlepage}%

	\begin{flushright}
	\includegraphics[height=1.0cm]{Figuras/logo_ESTG}
	\end{flushright}
	
	\vfill
	
  \begin{center}
  \Large{Trabalho Prático de }\\
  \Huge{\textbf{Matemática Discreta}}\\
  \Huge{\textbf{2024/2025}}
  	\end{center}  
  	
  	\vfill  	
  	\vfill
  	\vfill  	
  	\vfill
  	
  	Trabalho Elaborado por:\\
  	\textbf{Grupo M}
  	\begin{itemize}
  	    \item[] 
  	    8240231 David Sérgio Ferreira Alves
  	    \item []
  	    8240216 Pedro Afonso Farinha Gomes Lima Paraty
  	    \item[]
  	    8240266 Gabriel Alexandre Meireles Moreira
  	    \item[]
  	    8240839 Samuel Silva da Cunha
  	\end{itemize}
  	Curso de:\\
 \textsf{Licenciatura em Segurança Informática e Redes de Computadores}
	
  	\vfill  	
  	\vfill
  	
  	
  	
  	
  	Docentes:\\
  Eliana Costa e Silva (eos@estg.ipp.pt) \\
  Isabel Cristina Duarte (icd@estg.ipp.pt) 

	\begin{center}
  Felgueiras, \today
  	\end{center}  
        
\end{titlepage}%
\setcounter{footnote}{0}

%Página em branco - verso da Capa
\newpage\thispagestyle{empty}

%Índice
\tableofcontents
% Lista de Figuras
\listoffigures
% Lista de Tabelas
\listoftables

%Página em branco - verso do índice
\newpage\thispagestyle{empty}

%__________________________________________________________________________________________________________
% Pergunta 1
\chapter{Pergunta 1}

\subsection*{Conjuntos Escolhidos}

Seja o conjunto universo:
\[
U = \{1, 2, 3, \ldots, 20\}
\]

Escolheram-se os seguintes subconjuntos:
\[
A = \{3, 5, 9, 11, 13, 17, 19\} \quad (\#A = 7)
\]
\[
B = \{2, 3, 4, 5, 6, 7, 8, 9, 10, 11, 12, 13, 14, 15, 16, 17, 18\} \quad (\#B = 17)
\]

Estes conjuntos cumprem:
\begin{itemize}
    \item \(5 \leq \#A < 10\)
    \item \(\#B > 15\)
    \item \(A \neq B\), \(A \neq U\), \(B \neq U\)
\end{itemize}

\subsection*{Resolução das Alíneas}

\begin{enumerate}
    \item[a)] Cardinalidade:
    \[
    \#A = 7, \quad \#B = 17
    \]

    \item[b)] Complemento de B:
    \[
    \bar{B} = \{1, 19, 20\}
    \]
    \item[c)] União:
    \[
    A \cup B = \{2, 3, 4, 5, 6, 7, 8, 9, 10, 11, 12, 13, 14, 15, 16, 17, 18, 19\}
    \]

    \item[d)] Interseção:
    \[
    A \cap B = \{3, 5, 9, 11, 13, 17\}
    \]

    \item[e)] Diferença \(B - A\):
    \[
    B - A = \{2, 4, 6, 7, 8, 10, 12, 14, 15, 16, 18\}
    \]

    \item[f)] Diferença simétrica \(A \oplus B\):
    \[
    (A \setminus B) \cup (B \setminus A) = \{2, 4, 6, 7, 8, 10, 12, 14, 15, 16, 18, 19\}
    \]

    \item[g)] \(\overline{A \oplus {B}} \cup (A - B)\):

    Calculou-se primeiro o  $\overline{A \oplus B}$:

    \[
   \overline{A \oplus {B}} =U -  A \oplus {B} =\{1,3,5,9,11,13,17,20\}
    \]
De seguida:
    \[
    A - B = \{19\}
    \]
Por fim, realizou-se a reunião

    \[
    Resultado final: \{1,3,5,9,11,13,17,19,20\}
    \]

    \item[h)] Produto cartesiano \(B \times A\):

    Como existem no total 119 pares (17x7), selecionaram-se apenas os 5 primeiros pares:
    \[
    \{(2,3), (2,5), (2,9), (2,11), (2,13), \ldots\}
    \]

    \item[i)] Produto cartesiano \(A \times A \times A\):

      Como existem no total 343 trios (7$^{3}$), selecionaram-se apenas os 5 primeiros trios:
    \[
    \{(3,3,3), (3,3,5), (3,3,9), (3,3,11), (3,3,13), \ldots\}
    \]
\end{enumerate}



%__________________________________________________________________________________________________________
% Pergunta 2
\chapter{Pergunta 2}

Primeiro definimos o valor de $\beta$ como o último algarismo do número de estudante. Considerando que o número de estudante seja 8240266, temos que:

\[
\beta = 6
\]

Agora, escolhemos $n = 30$, de forma que satisfaça a condição $50 + \beta < 2n < 100 - \beta$. Ou seja, temos:

\[
50 + 6 < 2n < 100 - 6 \quad \Rightarrow \quad 56 < 2n < 94 \quad \Rightarrow \quad 28 < n < 47
\]



\begin{enumerate}


    \item[a)]
    Queremos calcular o somatório:
    \[
    \sum_{j=\beta+2}^{n} \left( \frac{-2\beta-1}{5} \right)^j
    \]
    Para \(\beta = 6\) e \(n = 30\), temos:
    \[
    \sum_{j= \beta+2}^{n} \left( \frac{-2\beta-1}{5} \right)^j = \sum_{j= 6+2}^{30} \left( \frac{-2 \cdot 6 - 1}{5} \right)^j
    \]
    Ou seja:
    \[
    \sum_{j= 8}^{30} \left( \frac{-13}{5} \right)^j
    \]


\section*{Resultado utilizando  código scilab:}
     \[
    \sum_{j= 8}^{30} \left( \frac{-13}{5} \right)^j = 2.032 \times 10^{12}
    \]
A sequência dos termos forma uma progressão geométrica com razão $\frac{-13}{5}$


   \item[b)]
   Queremos calcular o somatório:
   \[
\prod_{i \in C} \left( \frac{\beta + 1}{i } -1\right)^4, \quad C = \left\{ 5m \in \mathbb{Z}: m = 1, \dots, M \right\}, \quad M = \min \left( 5 + \beta, \left\lceil \frac{100}{\beta + 1} \right\rceil \right)
\]
    Para \(\beta = 6\) e \(n = 30\), temos:
     \[
\prod_{i \in C} \left( \frac{6 + 1}{i } -1\right)^4, \quad C = \left\{ 5m \in \mathbb{Z}: m = 1, \dots, M \right\}, \quad M = \min \left( 5 + \beta, \left\lceil \frac{100}{\beta + 1} \right\rceil \right) 
\]
 = \[
\prod_{i \in C} \left( \frac{7}{i }-1 \right)^4, \quad C = \left\{ 5m \in \mathbb{Z}: m = 1, \dots, M \right\}, \quad M = \min \left( 5 + \beta, \left\lceil \frac{100}{\beta + 1} \right\rceil \right) 
\]

\section*{Resultado utilizando  código scilab:}

\[
\prod_{i \in C} \left( \frac{7}{i } -1 \right)^4, \quad C = \left\{ 5m \in \mathbb{Z}: m = 1, \dots, M \right\}, \quad M = \min \left( 5 + \beta, \left\lceil \frac{100}{\beta + 1} \right\rceil \right) = 8.513^{-9}
\]

 \vspace{2em}
      \item[c)]

    Queremos calcular o somatório:

      \[
\prod_{k=1}^{n-15} \left( 3 \times \sum_{j=n-5}^{n} \left( \left\lfloor 1 + \frac{j + k}{200} \right\rfloor - \left\lceil \frac{6!}{\beta + 1} \right\rceil \right) \right)
\]

    Para \(\beta = 6\) e \(n = 30\), temos:
  \[
\prod_{k=1}^{30-15} \left( 3 \times \sum_{j=30-5}^{30} \left( \left\lfloor 1 + \frac{j + k}{200} \right\rfloor - \left\lceil \frac{6!}{6 + 1} \right\rceil \right) \right)
\]
 = \[
\prod_{k=1}^{15} \left( 3 \times \sum_{j=25}^{30} \left( \left\lfloor 1 + \frac{j + k}{200} \right\rfloor - \left\lceil \frac{6!}{7} \right\rceil \right) \right)
\]

\section*{Resultado utilizando  código scilab:}
\[
\prod_{k=1}^{15} \left( 3 \times \sum_{j=25}^{30} \left( \left\lfloor 1 + \frac{j + k}{200} \right\rfloor - \left\lceil \frac{6!}{7} \right\rceil \right) \right) = -9.080 \times 10^{48}
\] 

\end{enumerate}




\end{document}